% !TeX program = pdflatex
\documentclass[10pt,a4paper]{article}
\usepackage[top=14mm,bottom=17mm,left=17mm,right=17mm,footskip=8mm]{geometry}
\usepackage[T1]{fontenc}
\usepackage{lmodern,xcolor,enumitem,fancyhdr}
\usepackage[normalem]{ulem}
\usepackage[hidelinks,pdftitle={Jiajie Jin - Research CV},pdfauthor={Jiajie Jin}]{hyperref}
\definecolor{accent}{HTML}{B13C30}
\definecolor{ink}{HTML}{25272B}
\definecolor{muted}{HTML}{51555B}
\definecolor{rulecolor}{HTML}{A6A7A9}
\renewcommand{\familydefault}{\sfdefault}

\setlength{\parindent}{0pt}
\pagestyle{fancy}\fancyhf{}\renewcommand{\headrulewidth}{0pt}
\fancyfoot[R]{\footnotesize\color{muted}\thepage}
\newcommand{\blocktitle}[2]{\par\addvspace{8pt}\noindent{\fontsize{14}{17}\selectfont\bfseries\textcolor{accent}{#1}#2}\hspace{.65em}{\color{rulecolor}\leaders\hrule height 3.3pt depth -2.9pt\hfill\kern0pt}\par\nobreak\vspace{4pt}}
\newcommand{\entry}[4]{\par\addvspace{5pt}{\bfseries #1}\hfill{\small\color{accent}#2}\par{\small\color{muted}#3\hfill #4}\par}
\setlist[itemize]{leftmargin=1.35em,label=\textbullet,topsep=2pt,itemsep=1pt,parsep=0pt,partopsep=0pt}
\newcommand{\project}[3]{\par\addvspace{4pt}\href{#1}{\textbf{#2}}\enspace #3\par}
\newcommand{\me}{\uline{\textbf{J Jin}}}
\newcommand{\badge}[2]{\href{#1}{\textcolor{accent}{[#2]}}}
\newcommand{\pubitem}[5]{\par\addvspace{7pt}\begingroup\leftskip=2.4em\noindent\llap{\makebox[2.4em][l]{\small[#1]}}{\bfseries #2}\if\relax\detokenize{#3}\relax\else\enspace{\footnotesize #3}\fi\par{\small #4}\par{\small\color{muted}#5}\par\endgroup}

\begin{document}
\color{ink}\raggedright
\begin{center}
{\fontsize{26}{30}\selectfont\bfseries Jiajie Jin}\par\vspace{4pt}
{\color{accent}\bfseries Post-Training for Search\quad$\cdot$\quad Deep Research Agents\quad$\cdot$\quad Retrieval-Augmented Generation}\par\vspace{3pt}
{\small Beijing, China\quad|\quad\href{mailto:jinjiajie@ruc.edu.cn}{jinjiajie@ruc.edu.cn}\quad|\quad\href{https://github.com/ignorejjj}{github.com/ignorejjj}\quad|\quad\href{https://ignorejjj.github.io}{ignorejjj.github.io}}
\end{center}
I am a second-year PhD student at the Gaoling School of Artificial Intelligence, Renmin University of China, advised by Prof. Zhicheng Dou. I build autonomous agents that use information efficiently and produce results people can verify.\par\vspace{2pt}
I have published \textbf{20+ papers} at venues including NeurIPS, ICML, ICLR, ACL, and SIGIR. I am first or co-first author on \textbf{8 CCF-A papers}, including \textbf{5 oral presentations}. My work has \textbf{3,000+ Google Scholar citations} (h-index 21), with over 100 citations each for eight papers. Open-source projects I have led have earned \textbf{7,000+ GitHub stars} and \textbf{250,000+ Hugging Face downloads}.
\blocktitle{Research}{ Interests}
\begin{itemize}
\item \textbf{Autonomous agents.} Long-horizon planning and multi-agent coordination for human--agent collaboration, autonomous research, and iterative self-improvement (CoworkOS, Arbor, FlashRSI).
\item \textbf{Deep research.} Post-training foundation models for search, gathering evidence across sources, and tracing analytical conclusions back to that evidence (Atria-Dawn, Sage-2.1, WebThinker, HiRA, FinSight, VeriGraph).
\item \textbf{Retrieval-augmented LLMs.} Efficient retrieval and long-document refinement, with open-source tools for building and evaluating RAG systems (FlashRAG, LongRefiner, LaSER, BIDER).
\end{itemize}
\blocktitle{Edu}{cation}
\entry{Renmin University of China, Gaoling School of Artificial Intelligence}{Beijing}{Combined master's and PhD program in Artificial Intelligence, advised by Prof. Zhicheng Dou}{Sep 2023 -- Present}
\entry{University of Science and Technology of China}{Hefei}{BS in Statistics}{Sep 2019 -- Jun 2023}
\blocktitle{Research}{ Internships}
\entry{Shanghai AI Laboratory, Foundation LLM Team}{Shanghai}{Search Lead, mentored by Tao Gui}{Jul 2026 -- Present}
\begin{itemize}
\item Contributed to post-training Atria-Dawn and Sage-2.1. Led search task and trajectory synthesis, reinforcement learning optimization, and development of the agent harness for long-horizon tasks.
\item The team's models scored \textbf{84.6 on WideSearch}, \textbf{96.5 on DeepSearchQA}, and \textbf{91.3 on BrowseComp}. The service has reached 100,000 users, exceeded 2.78M calls in a day, and processed 400B tokens per day.
\end{itemize}
\entry{Microsoft Research Asia, Deep Search Team}{Beijing}{Research Intern, mentored by Kai Qiu and Chong Luo}{Mar -- Jun 2026}
\begin{itemize}
\item Led development of \href{https://github.com/RUC-NLPIR/Arbor}{\textbf{Arbor}}, an autonomous research agent that uses a hypothesis tree to track findings and revise experiment plans as results come in.
\item Arbor ranked first on held-out evaluations for all six research tasks, with \textbf{2.5$\times$ the average gain} of Codex and Claude Code. It achieved an \textbf{86.36\% Any Medal rate} on MLE-Bench Lite and earned 1.1k+ GitHub stars.
\end{itemize}
\entry{Alibaba Tongyi Lab, Search Agent Team}{Beijing}{Research Intern, mentored by Dingkun Long}{Sep 2025 -- Feb 2026}
\begin{itemize}
\item Developed \textbf{LaSER}, which embeds explicit reasoning in a dense retriever's latent space. It improved nDCG on BRIGHT by \textbf{5\%} and was adopted for internal retrieval tasks (SIGIR 2026, oral).
\end{itemize}
\blocktitle{Selected}{ Projects}
\project{https://github.com/RUC-NLPIR/FlashRAG}{FlashRAG}{Designed the modular RAG toolkit and independently implemented its core. It integrates 23 RAG methods, 7 information-agent methods, and 40 datasets. The project has \textbf{3.6k+ stars}, \textbf{250,000+ downloads} on Hugging Face, and 300+ paper citations.}
\project{https://arxiv.org/abs/2611.03178}{CoworkOS}{Led architecture and core development for multi-agent planning, execution, and review in shared Office, CAD, and Blender workspaces. Self-improvement raised the mean score across eight benchmarks from \textbf{33.19 to 42.95}.}
\project{https://github.com/RUC-NLPIR/FinSight}{FinSight}{A shared variable space connects financial data collection, report generation, chart validation, and evidence tracing. Won the \textbf{AFAC 2025 Challenge Track} (1st of 1,289) and an ACL 2026 SAC Highlight.}
\par\addvspace{3pt}{\small\textbf{Other open-source work}\enspace WebThinker (deep research, 1.5k+ stars) and FlashRSI (self-improvement, 420+ stars).}
\newpage
\blocktitle{Selected}{ Publications}
{\small\color{muted}First- or co-first-author papers. An asterisk denotes equal contribution.}\par
\pubitem{1}
    {FinSight: Towards Real-World Financial Deep Research}
    {\badge{https://github.com/RUC-NLPIR/FinSight}{Code}\,~\badge{https://arxiv.org/abs/2510.16844}{arXiv}}
    {\me, Y Zhang, Y Xu, H Qian, Y Zhu, Z Dou}
    {ACL 2026 Main, Oral, SAC Highlight}

\pubitem{2}
    {VeriGraph: Towards Verifiable Data Analysis Agents}
    {\badge{https://github.com/RUC-NLPIR/VeriGraph}{Code}\,~\badge{https://arxiv.org/abs/2606.16603}{arXiv}}
    {\me$^{*}$, Z Yang$^{*}$, W Liao, Y Hu, G Dong, X Li, Y Zhu, Z Dou}
    {NeurIPS 2026}

\pubitem{3}
    {LaSER: Internalizing Explicit Reasoning into Latent Space for Dense Retrieval}
    {\badge{https://github.com/RUC-NLPIR/LaSER}{Code}\,~\badge{https://arxiv.org/abs/2603.01425}{arXiv}}
    {\me, Y Zhang, M Li, D Long, P Xie, Y Zhu, Z Dou}
    {SIGIR 2026, Oral}

\pubitem{4}
    {HiRA: A Hierarchical Reasoning Framework for Decoupled Planning and Execution in Deep Search}
    {\badge{https://github.com/ignorejjj/HiRA}{Code}\,~\badge{https://arxiv.org/abs/2507.02652}{arXiv}}
    {\me, X Li, G Dong, Y Zhang, Y Zhu, Y Zhao, H Qian, Z Dou}
    {SIGIR 2026, Oral}

\pubitem{5}
    {Hierarchical Document Refinement for Long-Context Retrieval-Augmented Generation}
    {\badge{https://github.com/ignorejjj/LongRefiner}{Code}\,~\badge{https://arxiv.org/abs/2505.10413}{arXiv}}
    {\me, X Li, G Dong, Y Zhang, Y Zhu, Y Wu, Z Li, Q Ye, Z Dou}
    {ACL 2025 Main, Oral}

\pubitem{6}
    {FlashRAG: A Modular Toolkit for Efficient Retrieval-Augmented Generation Research}
    {\badge{https://github.com/RUC-NLPIR/FlashRAG}{Code}\,~\badge{https://arxiv.org/abs/2405.13576}{arXiv}}
    {\me, Y Zhu, Z Dou, G Dong, X Yang, C Zhang, T Zhao, Z Yang, JR Wen}
    {TheWebConf 2025 Resource Track, Oral}

\pubitem{7}
    {WebThinker: Empowering Large Reasoning Models with Deep Research Capability}
    {\badge{https://github.com/RUC-NLPIR/WebThinker}{Code}\,~\badge{https://arxiv.org/abs/2504.21776}{arXiv}}
    {X Li$^{*}$, \me$^{*}$, G Dong$^{*}$, H Qian, Y Zhu, Y Wu, JR Wen, Z Dou}
    {NeurIPS 2025, ranked \#1 on Hugging Face Daily Papers}

\pubitem{8}
    {BIDER: Bridging Knowledge Inconsistency for Efficient Retrieval-Augmented LLMs via Key Supporting Evidence}
    {\badge{https://arxiv.org/abs/2402.12174}{arXiv}}
    {\me, Y Zhu, Y Zhou, Z Dou}
    {ACL 2024 Findings}

\pubitem{9}
    {CoworkOS: Scaling Agent Coworkers for Real-world Automation}
    {\badge{https://github.com/RUC-NLPIR/CoworkOS}{Code}\,~\badge{https://arxiv.org/abs/2611.03178}{arXiv}}
    {\me, G Dong, C Pan, et al.}
    {arXiv:2611.03178, 2026 (preprint)}

\pubitem{10}
    {Toward Generalist Autonomous Research via Hypothesis-Tree Refinement (Arbor)}
    {\badge{https://github.com/RUC-NLPIR/Arbor}{Code}\,~\badge{https://arxiv.org/abs/2606.11926}{arXiv}}
    {\me$^{*}$, Y Hu$^{*}$, K Qiu, Q Dai, C Luo, G Dong, X Li, T Zhao, X Ma, G Zhang, Z Wu, B Liu, Z Yang, L Li, L Wang, H Qian, Y Zhu, Z Dou}
    {Preprint}


\blocktitle{Honors}{ and Awards}
\newcommand{\award}[2]{\noindent\makebox[1.4cm][l]{\small\color{muted}#1}#2\par\vspace{2pt}}
\award{2026}{Baidu Scholarship, global top 20}
\award{2026}{CAST Young Science and Technology Talent Cultivation Program}
\award{2026}{CIE--Tencent PhD Research Incentive Program, Hunyuan Track}
\award{2026}{BAAI Qingyuan InnoVibe Academic Rising Star (25 selected worldwide)}
\award{2026}{ACL SAC Highlight Paper Award for FinSight}
\award{2026}{Dean's Scholarship, Renmin University of China}
\award{2025}{First-Class Academic Scholarship, Renmin University of China}
\award{2025}{AFAC Financial Intelligence Challenge Track Champion (1st of 1,289)}
\award{2022}{JAC--NIO Special Scholarship, University of Science and Technology of China}
\award{2021}{First Prize, Anhui Division, National College Student Mathematics Competition}
\end{document}
